\lecture{6}{2. Marts 2026}{Physical interpretation of the strain tensor}

\subsection{Principal Strain}
As the strains have the same second-order tensor character as the stresses, it is expected that we are able to identify a system of principal strains that are the eigenvalues, determined from the characteristic equation
\begin{equation*}
  \left| \epsilon_{ji} - \lambda \delta_{ji} \right| = 0
\end{equation*}
which is analogous to the result for stresses.
The three roots $\lambda_1$, $\lambda_2$, and $\lambda_3$ are the principal strains $\epsilon^{(k)}$ with $k = 1,2,$ and 3.

The corresponding eigenvectors designate the directions associated with each of the principal strains and are computed from
\begin{equation*}
  \left( \epsilon_{ji} - \lambda \delta_{ji} \right)n_j = 0
.\end{equation*}
The directions of principal strain are denoted by $n_j^{(k)}$, with $k = 1, 2, $ and 3.
These directions are mutually perpendicular and, for isotropic materials, conincide with the directions of the principal stresses.
Which makes logical sense, as the largest stress is also expected to cause the largest strain.

\subsection{Volume and Shape Change}
Sometimes it is convenient to separate the components of strain into those that cause changes in volume and those that cause change in the shape of a differential element.

Considering a volume element oriented with the principal directions and with shearing stresses equal to zero.
The original length of each side is $\mathrm{d}x^{(k)}$ and the final length as $\mathrm{d}X^{(k)}$ where
\begin{equation*}
  \mathrm{d}X^{(k)} = \mathrm{d}x^{(k)} \left( 1 + \epsilon^{(k)} \right)
.\end{equation*}

The original volume is found as
\begin{equation*}
  \mathrm{d}v = \prod_{k = 1}^{3} \mathrm{d}x^{(k)}
\end{equation*}
and the final volume as
\begin{align*}
  \mathrm{d}V &= \prod_{k = 1}^{3} \mathrm{d}X^{(k)} \\
  &= \prod_{k = 1}^{3} \left( 1 + \epsilon^{(k)} \right) \, \mathrm{d}x^{(k)} \\
  &= \mathrm{d}v \prod_{k = 1}^{3} \left( 1 + \epsilon^{(k)} \right)
.\end{align*}

We define the \textit{dilation} as the relative change in volume
\begin{align*}
  \Delta &= \frac{\mathrm{d}V - \mathrm{d}v}{\mathrm{d}v} \\
  &= \frac{\mathrm{d}v \left( \left( \prod_{k = 1}^{3} \left( 1 + \epsilon^{(k)} \right) \right) - 1 \right)}{\mathrm{d}v} \\
  &= \left( \left( 1 + \epsilon^{(1)} \right)\left( 1 + \epsilon^{(2)} \right) \left( 1 + \epsilon^{(3)} \right) \right) - 1 \\
  &\approx \epsilon^{(1)} + \epsilon^{(2)} + \epsilon^{(3)} \\
  &= \prod_{k = 1}^{3} \epsilon^{(k)}
\end{align*}
which is the sum of the diagonal terms of the principal strain tensor.

The dilation $\Delta$ is also equal to the sum of the diagonal strains $\epsilon_{ii}$, and the divergence of the displacement vector $\nabla  \cdot \Vec{u} = u_{i,j}$.

Changes in shape of the volume element are called \textit{detrusions} and are denoted by shearing strains $\epsilon_{ij}, i \neq j$.

The separation of the strains into dilational and detrusional components can be formalized by resolving the strain tensor into two terms:
\begin{equation*}
  \Vec{\epsilon} = \Vec{\epsilon}_M + \Vec{\epsilon}_D
\end{equation*}
where
\begin{gather*}
  \Vec{\epsilon}_M = \begin{bmatrix}
    \epsilon_M & 0 & 0\\
  0 & \epsilon_M & 0\\
  0 & 0 & \epsilon_M\\
  \end{bmatrix} = \text{mean normal strain} \\
  \Vec{\epsilon}_D = \begin{bmatrix}
    \epsilon_{11} - \epsilon_M & \epsilon_{12} & \epsilon_{13}\\
  \epsilon_{12} & \epsilon_{22} - \epsilon_M & \epsilon_{23}\\
  \epsilon_{13} & \epsilon_{23} & \epsilon_{33} - \epsilon_M\\
  \end{bmatrix} = \text{strain deviator}
\end{gather*}
and where
\begin{align*}
  \epsilon_M &= \frac{1}{3} \left( \epsilon_{11} + \epsilon_{22} + \epsilon_{33} \right) \\
  &= \frac{1}{3} \left( \epsilon^{(1)} + \epsilon^{(2)} + \epsilon^{(3)} \right) \\
  &= \frac{1}{3} R_1
.\end{align*}

The mean normal strain corresponds to a state of equal elongation in all directions for an element at a given point.
Only subject to this, the element would maintain original shape but change in volume with dilation $\epsilon_{ii}$.
Hence, $\Vec{\epsilon}_M$ is the \textit{volumetric} component of strain.

The strain deviator, or \textit{deviatoric} component, $\Vec{\epsilon}_D$, has a dilation of
\begin{align*}
  \Delta &= \prod_{i = 1}^{3} \epsilon_D^{(i)} \\
  &= \epsilon_{D_{ii}} \\
  &= \left( \epsilon_{11}  -\epsilon_M \right) + \left( \epsilon_{22} - \epsilon_M \right) + \left( \epsilon_{33} - \epsilon_M \right) \\
  &= \epsilon_{11} + \epsilon_{22} + \epsilon_{33} - 3 \left( \frac{1}{3} \left( \epsilon_{11} + \epsilon_{22} + \epsilon_{33} \right) \right) \\
  &= 0
.\end{align*}
This corresponds to a change in shape with no change in volume of an element.

Corresponding to this, we can also define the principal stresses associated with changes in the volume and shape as
\begin{equation*}
  \Vec{\sigma} = \Vec{\sigma}_M + \Vec{\sigma}_D
\end{equation*}
defined similarly as the stress, where
\begin{gather*}
  \Vec{\sigma}_M = \begin{bmatrix}
    \sigma_M & 0 & 0\\
  0 & \sigma_M & 0\\
  0 & 0 & \sigma_M\\
  \end{bmatrix} = \text{mean normal stress} \\
  \Vec{\sigma}_D = \begin{bmatrix}
    \sigma_{11} - \sigma_M & \sigma_{12} & \sigma_{13}\\
  \sigma_{12} & \sigma_{22} - \sigma_M & \sigma_{23}\\
  \sigma_{13} & \sigma_{23} & \sigma_{33} - \sigma_M\\
  \end{bmatrix} = \text{stress deviator}
\end{gather*}
and where
\begin{align*}
  \sigma_M &= \frac{1}{3} \left( \sigma_{11} + \sigma_{22} + \sigma_{33} \right) \\
           &= \frac{1}{3} \left( \sigma^{(1)} + \sigma^{(2)} + \sigma^{(3)} \right)
.\end{align*}

In indicial form, the volumetric and deviatoric components of stress are, respectively,
\begin{align*}
  \sigma_{M_{ij}} &= \frac{1}{3} \sigma_{kk} \delta_{ij} \\
  \sigma_{D_{ij}} &= \sigma_{ij} - \sigma_{M_{ij}} \\
  &= \sigma_{ij} - \frac{1}{3} \sigma_{kk} \delta_{ij}
.\end{align*}

For principal axes, the stress deviator becomes
\begin{align*}
  \Vec{\sigma}_D &= \begin{bmatrix}
    \sigma^{(1)} - \sigma_M & 0 & 0\\
  0 & \sigma^{(2)} - \sigma_M & 0\\
  0 & 0 & \sigma^{(3)} - \sigma_M\\
  \end{bmatrix} \\
                 &= \begin{bmatrix}
                   \sigma_D^{(1)} & 0 & 0\\
                 0 & \sigma_{D}^{(2)} & 0\\
                 0 & 0 & \sigma_D^{(3)}\\
                 \end{bmatrix}
.\end{align*}
Worth noting here is that the sum of the diagonals of $\Vec{\sigma}_D$ is zero as
\begin{equation*}
  \sigma_{D_{ii}} = \sigma^{(1)} + \sigma^{(2)} + \sigma^{(3)} - 3 \cdot \frac{1}{3} \left( \sigma^{(1)} + \sigma^{(2)} + \sigma^{(3)} \right) = 0
.\end{equation*}

